\section{讲解笔记：【第3次作业笔记】 材料力学回顾(1)}
\begin{analysisbox}
本节按题目、答案和讲解笔记顺序整理。对原图中的结构几何关系，正文保留 可辨识文字，并在图形页中保留清理后的题图，便于核对杆件、支座、荷载和尺寸。
\end{analysisbox}
\subsection*{第 1 页转写}
结力基础班第3次作业

材料力学回顾

材料力学回顾作业及答案

、画剪力图和弯矩图

(a)

M图

ol

Q图

30kN

30 kv

15kN/m

1m

3m

(b)

ml= 30$\times$/ + $\times$15$\times$1= 375

30kN

15kN/m

315

1m

M (v.)

30

$\alpha$ (Iv )

\subsection*{第 2 页转写}
30kN-m

40kN

.5m

(c)

4kN m

0.2kN/m

Zmt=0

Me6 +0.4x / -1.4*2:0

Meto 2.4 fo.m

10m

1.4kn

mi\%: 1-4p2-+4 = 2.4

(d)

m图 (kv.m)

91 三,月$\times$70 x是

图(K\textasciitilde{})

\subsection*{第 3 页转写}
m

m

身

i

L

TukN

m uxl

zozkn.io

l

z

M图1kmmy

illQllknsii.li

in

it

T

fi

A

 B

qa

1

aj.lt

\subsection*{第 4 页转写}
二、试利用荷载集度、剪力和弯矩间的微分关系作下列各梁的剪力图和弯矩图。

5kN/m

1m

2.5m

a/ =Zxxs

(a)

M(kMM)

(m)n

15kN.10kN.m

1m

2m

(b)

MA: ID+ Xx/ = Vs

m(r.m)

$\alpha$ (km)

\subsection*{第 5 页转写}
3a

5a

1599

(c)

2.6594*

1.596

1.519

2M.

Me

3a

me

Me

(d)

ZMA=D

Fjox 54 + me- 2me: 0

me

2me

me

\subsection*{第 6 页转写}
4kN m

4kN/m

20kNm

4m

F9B=14kN

(e)

IMA-0 , 4FyB -4 - u4*2. -

(.m)

Q (kn)

4a

SFT

\subsection*{第 7 页转写}
2kN/m

(g）

M (w.n)

(KN)

250kN

250kN

0kN/m

  280kn

MD: Vgo$\times$-

180kn

(h)

20

M (lv.m )

40

260

()日

160

\subsection*{第 8 页转写}


qq2

悲

i

m 0

Fna

qaxz.co 可13 299



Tna l Bhai

g l T

ca

fo

o

MA 9axa qxax.ie i

廻

i

Fly

Dyuf

MA二0

问xG

99 9 0 Fg

qq.it

ii

ha

in

蜘

公1

00

结论：作 在铰上的集中力，可以认为是作 在铰的左边，也可以认为是作 在铰的右边。不影响最终的内力图。

\subsection*{第 9 页转写}
因边]

F=qa

Fuy=p

IMA=D,

Fay=0

M=qa

